\chapter{Fundamentals}
\label{ch:foundations}

This chapter introduces the forces on a moving appendage, two simple models of drag-based and lift-based
thrust, the metrics used to compare propulsors and the balance that turns single-leg thrust into a swimming
speed. The models are used to interpret the simulations, not to replace them.

\section{Forces on a Moving Appendage}
\label{sec:fund-forces}

A rigid appendage that moves through still water with velocity $\bm v$ relative to the water experiences a
pressure and shear force. Its components along and normal to $\bm v$ are the drag and the lift. For a
thin plate the drag depends strongly on the orientation: moved broadside, a plate has a drag coefficient of
about 1.2--2 \cite{Newman1977}; moved edgewise, it has only skin friction. Accelerating the plate also
accelerates the surrounding water, which appears as an added-mass force proportional to the acceleration.
Reduced-order models sum these contributions per body \cite{Morison1950,Fossen1994}:
\begin{equation}
  \bm F = -\tfrac12\rho\, C_D A\, |\bm v|\,\bm v \;+\; \bm L \;-\; C_a \rho V \dot{\bm v} \;+\; \bm F_\text{b},
  \label{eq:morison}
\end{equation}
with density $\rho$, reference area $A$, volume $V$, added-mass coefficient $C_a$ and buoyancy
$\bm F_\text{b}$. A pure drag model ($\bm L=\bm 0$) can only push water backwards; it is the model of a
\emph{paddle}. A lift term normal to $\bm v$ makes it possible to produce thrust with a foil that moves
mostly across the swimming direction; it is the model of a \emph{fin} or \emph{fluke}.

The flows of this thesis are unsteady and separated. With a paddle-tip speed of about
\SI{0.35}{\metre\per\second} and a fan radius of \SI{42}{\milli\metre}, the paddle works at a Reynolds number
$\mathit{Re}=vc/\nu\approx1.5\times10^4$. The fluke, with a chord of \SI{34}{\milli\metre} at
\SI{0.22}{\metre\per\second}, works at $\mathit{Re}\approx7.6\times10^3$. At these Reynolds numbers a starting
paddle forms a vortex whose growth sets the force \cite{Kim2011}, and a flapping foil carries a leading-edge
vortex \cite{Triantafyllou1991,Anderson1998}. Equation~\eqref{eq:morison} with constant coefficients misses both
effects. It is used here to derive trends, while magnitudes come from CFD.

\section{Quasi-Steady Paddling}
\label{sec:fund-paddle}

Consider a paddle that sweeps a path of length $L$ backwards (power stroke) with speed $v_p$ and area
$A$, and returns (recovery stroke) with speed $v_r=x\,v_p$ and a reduced area $aA$, $0<a\le 1$. At zero
forward speed the quasi-steady impulses of the two strokes are $\tfrac12\rho C_D A v_p L$ and
$-\tfrac12\rho C_D\,aA\,x v_p L$, and the period is $L(1+x)/(x v_p)$. The cycle-mean thrust is
\begin{equation}
  \Tm = \tfrac12\rho C_D A v_p^2\;\frac{x\,(1-a x)}{1+x}.
  \label{eq:paddle-mean}
\end{equation}
Both asymmetries identified by Herrera-Amaya and Byron \cite{HerreraAmaya2024} appear: a spatial one ($a<1$)
and a temporal one ($x\neq1$). For a fixed power-stroke speed, which is the natural constraint of a
speed-limited servo, \eqref{eq:paddle-mean} is maximal at
\begin{equation}
  x^* = \sqrt{1+1/a}-1, \qquad \mathrm{DUTY}^* = \frac{x^*}{1+x^*},
  \label{eq:duty-opt}
\end{equation}
where DUTY is the fraction of the cycle spent in the power stroke. A well-folding paddle (small $a$) should
return quickly, a poorly folding one slowly. For the muskrat, whose foot area is reduced by 55.5\,\% during
recovery ($a=0.445$), \eqref{eq:duty-opt} gives $\mathrm{DUTY}^*=0.445$, close to the measured 0.45
\cite{Fish1984}. For BODY2's fan paddle the effective ratio $a$ follows from the CFD
(\cref{sec:res-drag}); when $x^*>1$, the recovery would have to be faster than the power stroke, and the
servo speed limit then sets DUTY instead of \eqref{eq:duty-opt}.

Three further conclusions follow from the same model.
\begin{itemize}
  \item \emph{Velocity law.} For a given peak speed and stroke time, a trapezoidal velocity profile (short
        ramps, constant speed in between) produces a larger impulse $\int v^2\,\mathrm dt$ than a sinusoidal
        one. For a given impulse, H\"older's inequality shows that the energy $\int v^3\,\mathrm dt$ is
        minimal when $v$ is constant.
  \item \emph{Forward speed.} When the robot moves with speed $\U$, the power stroke acts on the relative
        velocity $v_p-\U$ and the recovery on $v_r+\U$. The thrust of the power stroke falls and the drag of
        the recovery rises with $\U$. The paddle should be open only while $v>\U$.
  \item \emph{Efficiency.} The useful power is $\Tm\U$, while the paddle delivers $\Tm v_p$ to the water during
        the power stroke. The Froude efficiency of a drag-based stroke is therefore bounded by
        $\U/v_p<1$ and is highest for a large paddle moved slowly \cite{Fish1996,Walker2000}.
\end{itemize}

\section{Flapping-Foil Kinematics}
\label{sec:fund-foil}

A foil that is carried with forward speed $\U$ and heaves with vertical position $h(t)$ meets the water at
the inflow angle
\begin{equation}
  \gamma(t) = \arctan\frac{\dot h(t)}{\U}.
  \label{eq:gamma}
\end{equation}
If the chord is pitched by $\theta(t)$ from the swimming direction, the effective angle of attack is
\begin{equation}
  \alpha(t) = \gamma(t) - \theta(t),
  \label{eq:alpha}
\end{equation}
with signs chosen so that $\theta$ and $\gamma$ rotate the same way (\cref{fig:foil-sketch}). The lift $L$
acts normal to the relative inflow and the drag $D$ along it; the thrust is
\begin{equation}
  F_x = L\sin\gamma - D\cos\gamma .
  \label{eq:foil-thrust}
\end{equation}
Thrust is produced as long as the lift is tilted forwards enough to overcome the drag. Because $\sin\gamma$
multiplies the lift, the same foil at the same $\alpha$ produces a forward force whenever it heaves; unlike the
paddle, it does not need a recovery stroke against the flow.

\begin{figure}[t]
  \centering
  \begin{tikzpicture}[scale=1.1,>=Latex]
    \draw[->,gray] (1.8,-2.0) -- (-1.2,-2.0) node[left,font=\small,black]{swimming direction $x$};
    \draw[gray,dashed] (-3.4,0) -- (2.2,0);
    \begin{scope}[rotate=-15]
      \fill[TUMblue!75] (-1.3,0) .. controls (-1.2,0.2) and (-0.4,0.16) .. (1.3,0)
                        .. controls (-0.4,-0.12) and (-1.2,-0.18) .. (-1.3,0);
      \draw[TUMblue!40!black,dotted] (-2.2,0) -- (2.0,0);
    \end{scope}
    \draw[->] (1.9,0) arc (0:-15:1.9); \node[font=\small] at (2.15,-0.28) {$\theta$};
    \draw[->,thick] (-3.3,1.87) -- (-1.55,0.40) node[pos=0.25,above right,font=\small]{relative inflow};
    \draw[->] (-2.35,0) arc (180:140:0.9); \node[font=\small] at (-2.55,0.35) {$\gamma$};
    \draw[->,thick,TUMblue] (0.9,0.35) -- (0.9,1.25) node[above,font=\small]{heave $\dot h$};
    \draw[->,very thick,todo] (0,0) -- (-0.84,-1.0) node[below left,font=\small]{$L$};
    \draw[->,very thick,todo!55] (0,0) -- (0.54,-0.45) node[right,font=\small]{$D$};
  \end{tikzpicture}
  \caption[Flapping-foil angles]{Angles of a heaving and pitching foil. The relative inflow is tilted by the
    inflow angle $\gamma=\arctan(\dot h/\U)$ below the swimming direction; the chord is pitched by $\theta$
    in the same sense; the angle of attack between chord and inflow is $\alpha=\gamma-\theta$. Lift $L$ acts
    normal and drag $D$ parallel to the relative inflow; the thrust is their net component along $x$. Shown
    for the upstroke.}
  \label{fig:foil-sketch}
\end{figure}

The Strouhal number $\mathit{St}=fA/\U$, with the peak-to-peak excursion $A$ of the trailing edge, is the ratio of heave speed to forward speed. In this thesis $A$ is taken as the excursion of the pin, which is known from the prescribed motion; the trailing-edge value is somewhat larger. Efficient cruising lies at $0.2<\mathit{St}<0.4$
\cite{Taylor2003,Triantafyllou1991}. A leg driven by a speed-limited servo has an approximately fixed heave
speed $\dot h_\text{max}$. The Strouhal number, and with it the inflow angle, then decreases with swimming speed:
\begin{equation}
  \gamma_\text{max}(\U) = \arctan\frac{\dot h_\text{max}}{\U}.
  \label{eq:gammamax}
\end{equation}
If the pitch schedule $\theta(t)$ is kept fixed while $\U$ changes, the angle of attack changes as well. At rest
($\U\to0$) the inflow is vertical, $\gamma=\SI{90}{\degree}$, and the fluke is stalled; it works as a small
paddle. At high speed $\gamma_\text{max}$ becomes small and a fixed pitch amplitude can make $\alpha$ too small
or even negative.

A simple way to keep $\alpha$ in the efficient range of \SIrange{15}{25}{\degree} \cite{Read2003,Hover2004} is to
scale the pitch with the inflow angle,
\begin{equation}
  \theta(t) = \theta_0\,\frac{\dot h(t)}{\dot h_\text{max}}, \qquad \theta_0=\tfrac12\,\gamma_\text{max}(\U),
  \label{eq:l3}
\end{equation}
so that at mid-stroke, where $\dot h=\dot h_\text{max}$, the foil takes half of the inflow angle as pitch and
the other half as angle of attack. This \emph{speed-scheduled pitch} is evaluated in
\cref{sec:res-l3}. It requires one parameter per speed and no knowledge of the flow. For BODY2's leg,
$\dot h_\text{max}\approx\SI{0.28}{\metre\per\second}$, so $\gamma_\text{max}=\SI{52}{\degree}$,
\SI{43}{\degree} and \SI{35}{\degree} at \SIlist{0.223;0.30;0.40}{\metre\per\second}.

\section{Performance Metrics}
\label{sec:fund-metrics}

All single-leg quantities are cycle means over complete stroke periods $T=1/f$. For a propulsor with
position $\bm x_H(t)$ of its reference point, angular velocity $\bm\omega(t)$ and fluid force and moment
$\bm F$, $\bm M_H$ about that point:
\begin{align}
  \Tm &= \big\langle F_x \big\rangle, \label{eq:thrust}\\
  \Pm &= -\big\langle \bm F\cdot\dot{\bm x}_H + \bm M_H\cdot\bm\omega \big\rangle, \label{eq:power}\\
  \etaF &= \frac{\Tm\,\U}{\Pm}, \qquad C_T=\frac{\Tm}{\tfrac12\rho\U^2S}. \label{eq:eta}
\end{align}
The $x$ axis points in the swimming direction, so that $\Tm>0$ is thrust and $\Pm>0$ is power delivered
to the fluid. $\Pm$ is the hydrodynamic input power: it does not include the inertia of the parts, friction
or motor losses. $\etaF$ is the Froude (propulsive) efficiency; it vanishes at $\U=0$, where the thrust per
power $\Tm/\Pm$ is reported instead. $S$ is the planform area of the propulsor. For the whole robot the cost of
transport is reported either as energy per distance $\Pm/\U$ (J/m) or dimensionless as $\Pm/(mg\U)$; the
form used is stated with each value.

For the drag-type paddle, the thrust is also split into the contributions of the power stroke and the
recovery,
\begin{equation}
  \Tm = \Tm_\text{pow} + \Tm_\text{rec}, \qquad
  \Tm_\text{pow} = \frac1T\int_{\text{power}} F_x\,\mathrm dt, \quad
  \Tm_\text{rec} = \frac1T\int_{\text{recovery}} F_x\,\mathrm dt,
  \label{eq:split}
\end{equation}
which shows directly which half of the cycle limits the thrust at speed.

\section{Self-Propulsion Balance}
\label{sec:fund-selfprop}

A robot swims steadily at the speed where the propulsive force of its legs equals the resistance of the
hull and the parts that do not propel. If each of the four legs produces the cycle-mean thrust $\Tm(\U)$
measured in isolation and the legs do not interact, the cruising speed $\U_c$ solves
\begin{equation}
  4\,\Tm(\U_c) = D(\U_c),
  \label{eq:selfprop}
\end{equation}
and the time history from rest follows from the quasi-steady surge equation
\begin{equation}
  (m+m_a)\,\dot\U = 4\,\Tm(\U) - D(\U).
  \label{eq:surge}
\end{equation}
Equation~\eqref{eq:selfprop} neglects the interaction between the legs, which can raise the thrust of four
paddling legs above four times the single-leg value \cite{Wang2025}, and the effect of the leg motion on the
hull resistance. It also assumes that the legs keep their prescribed motion at every speed. Because the
resistance of BODY2 is known only within bounds, \eqref{eq:selfprop} is solved for three resistance
curves in \cref{sec:res-selfprop}, and the result is reported as a range.
